\documentclass[11pt]{article}
\usepackage{fontspec}
\usepackage[french,english]{babel}
\usepackage{amsmath}
\usepackage{amssymb}
\usepackage{amsthm}
\usepackage{mathrsfs}
\usepackage{tikz-cd}
\usepackage{enumitem}

\newcommand{\printedpage}[1]{\marginpar{\footnotesize p.~#1}}
\newcommand{\msp}[1]{\marginpar{\footnotesize [#1]}}
\newcommand{\Ch}{\underline{\mathrm{Ch}}}
\newcommand{\CCh}{\underline{\underline{\mathrm{Ch}}}}
\newcommand{\Xb}{\overline{X'}}
\newcommand{\Xbb}{\overline{\overline{X'}}}
\newcommand{\Hotz}{\mathrm{Hot}^{\wedge}_{0}}
\newcommand{\Ty}{\underline{\mathrm{Ty}}}

\begin{document}

\printedpage{94}
% head: A. Grothendieck : Dérivateurs, Ch. VII, Édition des Universités de Montpellier II et Paris VII
retrograde, $\Ch_{n,m}(X) \simeq \Ch_{\tau_{n,m}}(X)$ is Cat-\textit{fibered} over $X \times X$, hence smooth (i.e. HOT-smooth), and this passes to $\varinjlim$ (\footnote[36]{\textbf{NB} $\Ch_\infty(X)$ is neither fibered nor cofibered over $X \times X$ - for the inclusion functors between the $\Ch_{n,m}$ are not cartesian resp. cocartesian.}).

Thus $\Ch_\infty(X) \xrightarrow{p} X \times X$ is both proper and smooth (\footnote[37]{12.11. But that already implies that it is a fibration!}). To say that it is a HOT-fibration, amounts e.g. to saying that for
\[
(a,b) \xrightarrow{(u,v)} (a',b')
\]
in $X \times X$, the corresponding arrow
\[
\Ch_\infty(X)/(a,b) \longrightarrow \Ch_\infty(X)/(a',b')
\]
is a Hot-equivalence. Now, since $p$ is proper, the inclusions of the fibers at $(a,b)$ and at $(a',b')$ into the respective terms are Hot-equivalences:
\msp{page 151}
\[
\begin{tikzcd}[row sep=large]
\Ch_\infty(X)/(a,b) \arrow[r, "\rho"] & \Ch_\infty(X)/(a',b') \\
\Ch_\infty(X;a,b) \arrow[u, hook, "\wr"] \arrow[r, dotted, "\varphi"] & \Ch_\infty(X;a',b') \arrow[u, hook, "\wr"]
\end{tikzcd}
\tag*{$(*)$}
\]
Thus we find a dotted arrow in Hot, which we have to prove is a Hot-equivalence. But by considering
\[
u : a \longrightarrow a' \qquad v : b \longrightarrow b'
\]
as \textit{paths} in $X$, and
\[
u^o : a' \longleftarrow a
\]
as a retrograde path from $a'$ to $a$, the general variance theory (in IX) does indeed give us an isomorphism in (Hot):
\[
\Ch_\infty(X;a,b) \longrightarrow \Ch_\infty(X;a',b'),
\]
we feel like saying that it is $\varphi$, i.e. that by inserting it in place of $\varphi$, the square above becomes commutative (which would prove what we want). I think that if we work with the ind-objects $\CCh_\infty(X)/(a,b)$, $\CCh_\infty(X;a,b)$ etc., for which we have pairings which can be spelled out nicely at the finite stages, and give rise to the $\varphi$, this should be easy (though a little tedious). This would not for all that prove the commutativity of $(*)$ in Hot, but
% foot: Version 22 décembre 2001
\printedpage{95}
% head: A. Grothendieck : Dérivateurs, Ch. VII, Édition des Universités de Montpellier II et Paris VII
\msp{page 152}
only in $\Hotz$, which is perhaps not sufficient. But this suffices to prove that the arrow $\rho$ in $(*)$ becomes \texttt{[an]} isomorphism in $\Hotz$, hence it is \texttt{[an]} isomorphism.

2$^o$) We have to show that the HOT-fibrations
\[
\Ch_\infty(X) \overset{\sigma}{\underset{\beta}{\rightrightarrows}} X
\]
are trivial, or what amounts to the same thing, that they have aspheric fibers. It suffices to see this for $\sigma$, hence that
\[
\Ch_\infty(X;a,?)
\]
is aspheric. But this category is the direct limit of the $\Ch_n(X;a,?)$, i.e. of the
\[
\underline{\mathrm{Hom}}^{\text{fixed source}}(\tau^n, X).
\]
But since the interval $I_{\tau^n}$ retracts by deformation onto its origin, it is immediate that this \underline{Hom} category is even contractible.

3$^o$) Let us consider
\[
\begin{tikzcd}[column sep=small, row sep=large]
 & \Xb = C(f) \arrow[dl] \arrow[dr, "g"] \arrow[ddrr, bend left, "\bar f"] & & \\
X' \arrow[ur, dotted, bend left, "\alpha"] \arrow[dr, "f"'] & \underline{\text{cart.}} & \Ch_\infty(X) \arrow[dl, "\sigma"'] \arrow[dr, "\beta"] & \\
 & X & & X,
\end{tikzcd}
\]
\msp{page 153}
I would like to prove that $\bar f$ is a proper and smooth \textbf{D}-fibration (\footnote[38]{That follows formally from what has already been seen, cf. pages 170-171.}).

$\Xb$ is the category of pairs
\[
(x',c) \ \in \ \mathrm{Ob}\,X' \times \mathrm{Ch}_\infty(X)
\]
such that
\[
\sigma(c) = f(x')
\]
and
\[
\bar f(x',c) = \beta(c), \qquad \text{target of } c.
\]
The fiber of $\Xb$, $\bar f$ at $x_0 \in X$ is thus the category of pairs $(x',c) \in \mathrm{Ob}\,X' \times \mathrm{Ch}_\infty(X)$, such that
\[
\sigma(c) = f(x'), \qquad \beta(c) = x_0.
\]
% foot: Version 22 décembre 2001

\printedpage{96}
% head: A. Grothendieck : Dérivateurs, Ch. VII, Édition des Universités de Montpellier II et Paris VII
We can also regard $\Xb$ as the direct limit of the subcategories $\Xb_{n,m}$, corresponding to the pairs $(x',c)$ with $c \in \mathrm{Ch}_{n,m}(X)$. We find that for $m$ even (terminal direction of $\tau_{n,m}$ retrograde) the $\Xb_{n,m}$ are Cat-fibered over $X$, for $m$ odd (terminal direction of $\tau_{n,m}$ direct) they are Cat-cofibered. So we conclude again as before, $\Xb$ being (at will) a filtering direct limit of subcategories smooth over $X$, or proper over $X$. % [sic] correspondante -> correspondantes, agreement with « sous-catégories »

\msp{page 154}
Finally we have to prove that $\Xb \xrightarrow{\bar f} X$ is a \textbf{D}-fibration (\footnote[39]{This already follows from the fact that it is proper and smooth!}), which amounts e.g. to saying that for an arrow
\[
a \xrightarrow{v} b
\]
in $X$, the induced arrow
\[
\Xb/a \xrightarrow{\bar v} \Xb/b
\]
is a Hot-equivalence. We already know as before that we have a diagram with Hot-equivalences as vertical arrows
\[
\begin{tikzcd}[row sep=large]
\Xb/a \arrow[r] & \Xb/b \\
\Xb_a \arrow[u, hook, "\wr"] \arrow[r, dotted] & \Xb_b, \arrow[u, hook, "\wr"]
\end{tikzcd}
\]
and it all comes down to showing that a certain arrow between $\Xb_a$ and $\Xb_b$ in Hot is \texttt{[an]} isomorphism. I feel like defining such an arrow directly in Hot, or even in Cat. Now for odd $m$, the $\Xb_{n,m}$ are Cat-cofibered over $X$. For $\Xb_{-m,m}$ ($m$ odd) we therefore find
\[
(\Xb_{-m,m})_a \xrightarrow{v_*} (\Xb_{-m,m})_b,
\]
\msp{page 155}
the trouble is that they are not compatible with the transition morphisms for variable $m$. They surely are, however, \textit{up to homotopy}, so we will find at least an arrow of ind-objects in Cat\,hot
\[
\Xbb_a \xrightarrow{v_*} \Xbb_b,
\]
which by passing to $\Hotz$ gives an arrow all the same
\[
\Xb_a \xrightarrow{v_*} \Xb_b.
\]
Working with odd $m$, we likewise find functors in the opposite directions % [sic] m impairs -> m pairs, the Cat-fibrées (giving v^*) are stated above for m pair
\[
(\Xb_{-m,m})_a \xleftarrow{v^*} (\Xb_{-m,m})_b,
\]
% foot: Version 22 décembre 2001
\printedpage{97}
% head: A. Grothendieck : Dérivateurs, Ch. VII, Édition des Universités de Montpellier II et Paris VII
which do not glue together for variable $m$, but which glue together up to homotopy, to give an arrow in Ind\,Cat\,hot
\[
\Xbb_a \xleftarrow{v^*} \Xbb_b,
\]
whence also
\[
\Xb_a \xleftarrow{v^*} \Xb_b \qquad \text{in } \Hotz.
\]
On the other hand, we should be able to show soritally,
\msp{page 156}
stage by stage, that $v_*v^*$ and $v^*v_*$ in Ind\,Cat\,hot are equal to the identity, hence $v_*$ and $v^*$ are isomorphisms inverse to each other. We must check on the other hand soritally that
\[
\begin{tikzcd}[row sep=large]
\Xbb/a \arrow[r, "\bar v"] & \Xbb/b \\
\Xbb_a \arrow[u] \arrow[r, "v_*"] & \Xbb_b, \arrow[u]
\end{tikzcd}
\]
in Ind\,Cat\,hot is commutative. Passing to the corresponding diagram in $\Hotz$, by taking direct limits, we find a commutative square in $\Hotz$
\[
\begin{tikzcd}[row sep=large]
\Xb/a \arrow[r, "\bar v"] & \Xb/b \\
\Xb_a \arrow[u, "\wr"'] \arrow[r, "v_*", "\sim"'] & \Xb_b, \arrow[u, "\wr"']
\end{tikzcd}
\]
whence it follows that $\bar v$ is \texttt{[an]} isomorphism in $\Hotz$, hence also in Hot.

To check that the HOT-fibration $\bar f$ is trivial if and only if $f$ is a HOT-equivalence, we would have to have developed in Cat the formalism of the exact sequence of loops - we are very close, evidently!
% foot: Version 22 décembre 2001

\printedpage{98}
% head: A. Grothendieck : Dérivateurs, Ch. VII, Édition des Universités de Montpellier II et Paris VII
\msp{page 157}

\begin{flushright}
\footnotesize 13.11.90
\end{flushright}

\section*{13\quad Where we finally prove that $\Ch(X;-,-)$ is as good as (or even better than) $\Ch_\infty(X;-,-)$.}

I believe after all that the categories $\Ch(X)$, $\Ch(X;a,b)$ (without the subscript $\infty$!) are ``good'', that the canonical arrows in Hot
\[
\left\{
\begin{aligned}
\Ch_\infty(X) \quad &\xrightarrow{\sim} \quad \Ch(X)\\
\Ch_\infty(X;a,b) \quad &\xrightarrow{\sim} \quad \Ch(X;a,b)
\end{aligned}
\right.
\]
are isomorphisms. (The first is even \texttt{[an]} isomorphism in $\mathrm{HOT}(X \times X)$, which implies the second relation, on passing to the fibers \dots) I believe that I am now (after XI) in a position to prove it. This will follow from the following two key results:

\begin{enumerate}
\item[1)] The canonical functor
\[
X \longrightarrow \Ch(X)
\]
(associating to every $x \in \mathrm{Ob}\,X$ the empty path with source and target $X$) is a Hot-equivalence, and: % [sic] but $X$ -> but $x$, the object is $x \in \mathrm{Ob}\,X$
\item[2)] The source-target functor
\[
\Ch(X) \longrightarrow X \times X
\]
is \textit{smooth} and \textit{proper} (hence a ``perfect Hot-fibration'').
\end{enumerate}

This should give what we are looking for. Since the $X \times X$-arrow $\Ch_\infty(X) \longrightarrow \Ch(X)$ (and the deduced arrows $\Ch_\infty(X;a,b) \longrightarrow \Ch(X;a,b)$) was defined by way of the category
\[
\mathcal{C}_{\mathbf{N}^o} = \mathbf{N}^o \times_{\Ty} \mathcal{C}
\qquad
\begin{aligned}
&\text{(where $\mathcal{C} = \Ch(X)$ \texttt{[is]} considered as a fibered category over $\Ty$,}\\
&\text{and $\mathbf{N}^o \longrightarrow \Ty$ is the projective system of the $\underbrace{(\Delta_+\Delta_-)^{2n}}_{t_n} \dots$)}
\end{aligned}
\]
\msp{page 158}
with the help of the two canonical arrows
\[
\begin{tikzcd}[row sep=huge, column sep=large]
\mathcal{C}_{\mathbf{N}^o} \arrow[r] \arrow[d] & \Ch_\infty(X) \\
\Ch(X), &
\end{tikzcd}
\]
% foot: Version 22 décembre 2001

\end{document}
